\documentclass[book.tex]{subfiles}

\begin{document}
\chapter{Group Actions and Representations}

\label{chap:lorentz_group}
\noindent\rule{11cm}{0.4pt} \\
\textbf{Prerequisites:} \autoref{chap:lie} \\
\textbf{Difficulty Level:} ** \\
\noindent\rule{11cm}{0.4pt}
\vspace{5mm}

Mathematical groups represent collections of mathematical transformations (such as rotations, translations, reflections) abstract. Representations provide a tool to turn abstract transformations into concrete matrices. In this chapter, we define the useful notions of group actions and group representations, which provide tools to study groups as concrete matrices instead of as abstract mathematical entities.

\section{Group Action}

A group $G$ is said to act on space $A$ if there exists some function 
\begin{align}
G \times A \to A
\end{align}
This function is commonly indicated by $\cdot$ under the notation
\begin{align}
g \cdot a \in A
\end{align}
Group representations furnish a natural group action on their associated vector spaces $V$.

\section{Defining Group Representations}

For example, a group representation on $\mathbb{R}^n$ is given by a map
\begin{align}
\rho: G \to M_{n\times n}(\mathbb{R})
\end{align}
We say that $\rho$ must be a group homomorphism that satisfies the following property
\begin{align}
    \rho(g_1g_2) = \rho(g_1)\rho(g_2)
\end{align}
for all $g_1, g_2 \in G$.

A faithful representation is one which is injective. We can generalize this notation of representation to the complex numbers
\begin{align}
\rho: G \to M_{n \times n}(\mathbb{C})
\end{align}
or even further to an arbitrary vector space $V$
\begin{align}
\rho: G \to M_{n \times n}(V)
\end{align}

To provide a concrete example, consider the group of 3 elements, 
\begin{align}
H &= \{1, u, u^2\}
\end{align}
where
\begin{align}
u &= e^{2\pi i / 3}
\end{align}
We define $\rho: H \to M_{2 \times 2}(\mathbb{C})$ by
\begin{align}
\rho(0) &= \begin{pmatrix}
1 & 0 \\
0 & 1 
\end{pmatrix}\\
\rho(u) &= \begin{pmatrix}
1 & 0 \\
0 & u 
\end{pmatrix} \\
\rho(u^2) &= \begin{pmatrix}
1 & 0 \\
0 & u^2 
\end{pmatrix} 
\end{align}
You will show in the exercises that $\rho$ is a group representation.


\section{Irreducible Representations}
Suppose that $W$ is a vector subspace of $V$ which is invariant under the group representation $\rho$. That is, if $w \in W$, we have that for all group elements $g \in G$,
\begin{align}
\rho(g) w \in W
\end{align}
That is, the matrices $\rho(g)$ map elements of $W$ back into $W$. We can then say that $W$ provides a subrepresentation of $\rho$. We say that $V$ provides an irreducible representation if there do not exist any subrepresentations. Irreducible representations are very useful since reducible representations can be decomposed into a sum of irreducible representations. For this reason, irreducible representations furnish the "atomic" elements of representations. In fact, there is a deep connection between particles and representation theory given by Wigner's theorem.



\section{Exercises}
\begin{enumerate}
    \item Prove that $H$ is a group. 
    \item Prove that $\rho$ acting on $H$ defined above is a group representation.
\end{enumerate}

\end{document}